\documentclass[12pt]{article}
\usepackage[margin=1in]{geometry}
\usepackage[T1]{fontenc}
\usepackage[utf8]{inputenc}
\usepackage{lmodern}
\usepackage{microtype}
\usepackage{amsmath,amssymb,mathtools}
\usepackage{enumitem}
\usepackage{setspace}
\usepackage{xurl}
\usepackage[hidelinks]{hyperref}
\setstretch{1.60}
\setlength{\parindent}{0pt}
\setlength{\parskip}{6pt plus 2pt minus 1pt}
\setcounter{secnumdepth}{-\maxdimen}
\emergencystretch=3em
\sloppy
\begin{document}
\section{1--2. Complete systems of invariants for the ternary biquadratic form}

\textit{Original publications: preliminary report, 1907; dissertation article, 1908.}

The problem is the construction of the complete system of concomitants belonging to the ternary
quartic form. I shall use the symbolic method, because it keeps visible the two operations that are
decisive: polarization and contraction. Let \[   f=a_x^4=b_x^4=c_x^4=\cdots \] be the symbolic
notation for the ground form. The symbolic letters are not independent coefficients, but devices for
writing invariant differential operations. An expression built from determinants such as \[
(abc)=\det(a,b,c) \] and from factors such as \(a_x,b_x,c_x\) is to be read as a covariant after
umbral substitution. The task is therefore not merely to list expressions, but to determine a finite
generating set from which every symbolic expression of the correct total degree follows by algebraic
combination.


The ternary biquadratic is the first case in which the older methods of transvectants cease to be
short. The number of possible symbolic products grows rapidly, and many apparent forms are
reducible. One must distinguish at every stage between a newly appearing form and a form obtained
from lower ones by multiplication, by applying the Aronhold operator, or by taking a polar. The
guiding principle is that every concomitant can be assigned a type: its degree in the coefficients,
its order in the variables, and its class in the dual variables. Forms of different type cannot
cancel. This permits the computation to be organized by tables rather than by trial.


I use the following elementary reductions repeatedly. If a symbolic product contains two equal
symbolic letters in a determinant, the product vanishes. If the same determinant occurs with its
letters permuted, only the sign changes. If a symbolic letter occurs to a power higher than the
degree allowed by the ground form, the expression can be replaced by a linear combination with lower
powers. Finally the fundamental identities among determinants in three letters, obtained by
expanding a four-by-four determinant with two identical rows, give all linear syzygies needed for
the reduction of the symbolic monomials.


For example, an expression of the form \[   (abc)^2 a_x^2 b_x^2 c_x^2 \] is not treated as an
isolated object. It is compared with all products obtained by replacing one occurrence of \(x\) by a
polar variable, by contracting two letters, and by multiplying invariants of smaller degree. If
after all these operations no reduction is possible, the expression is entered as a member of the
fundamental system. If a reduction is possible, it is recorded in the table of dependencies. The
same procedure is then repeated at the next symbolic degree.


The computation naturally divides into invariants, covariants, contravariants, and mixed
concomitants. The invariants have no remaining variables. The covariants contain only point
variables; the contravariants contain only line variables; the mixed forms contain both. The passage
from covariants to contravariants is performed by the dual operation, but one must not suppose that
this simply halves the work. Some forms are self-dual only after multiplication by a determinant
factor, and some reductions have no visible dual until they are written in normalized symbolic form.


The Aronhold process is used in the following way. Suppose that a covariant of order \(r\) has been
found. Applying the invariant differential operator attached to a second symbolic copy of the ground
form produces a new covariant of order \(r+4-2s\), where \(s\) is the number of contractions
performed. Terms that contain a factor already known to be reducible are immediately discarded. The
remaining terms are compared against the current table. In this manner all forms of a given degree
are obtained from forms of smaller degree together with a finite list of primitive exceptions.


The result of the calculation is a finite complete system. It is important that the word complete be
understood in the strict sense: every invariant and every concomitant of the ternary quartic is an
integral rational polynomial in the listed forms. There is no appeal to analytic continuity or to a
generic specialization. The proof consists of the symbolic reduction itself, together with the
verification that the reductions are compatible with specialization of the coefficients. Thus the
system is not a system only for the general quartic, but for every ternary quartic.


The reductions can be arranged according to weight. Let \(M_{d,r,s}\) denote the module of symbolic
concomitants of coefficient degree \(d\), point order \(r\), and line order \(s\). The products of
earlier members span a submodule \(P_{d,r,s}\). What has to be found is a basis of a complement of
\(P_{d,r,s}\) in \(M_{d,r,s}\). The symbolic identities give a presentation of \(M_{d,r,s}\); hence
the computation is finite. In the cases where the quotient is zero, no new form occurs. Where it is
non-zero, the number of new fundamental forms is exactly its rank.


This also explains why the construction is independent of the symbolic notation. If one writes the
quartic as a polynomial in coefficients \(a_{ijk}\), the concomitants become polynomials in these
coefficients and in the auxiliary variables. The symbolic calculus is simply a compressed notation
for these polynomials. Every reduction identity is an identity in the coefficient ring. Thus the
final system is a system in the ordinary algebraic sense.


The decisive point of the dissertation is therefore not one particular invariant, but the method of
closing the table. After a certain degree the Aronhold operations and products account for every
symbolic type. This is checked by the dimension count supplied by the Cayley series. The coefficient
of the corresponding term in the series gives the number of linearly independent concomitants of
that type. The table of already generated products reaches this number; consequently there are no
missing primitive forms.


In this form the computation is also a test of the general finiteness philosophy in invariant
theory. The individual formulae are long, and their printed arrangement is necessarily compressed.
But the proof does not rely on a guess that the list is complete. Completeness follows from the
comparison between the symbolic reduction and the enumerative series. The list is therefore not a
catalogue in the empirical sense, but the outcome of a finite algorithm.


I shall use in the sequel the following convention for the tables. A new form is marked by its
symbolic representative, followed by its degree and its order-class pair. If two representatives
differ only by a determinant identity, only one is retained. If a form is obtained by
multiplication, it is written in the product column and not repeated as a primitive form. If a form
arises from a transvection of lower forms, the generating operation is recorded. This convention
makes the tables reproducible and prevents the same concomitant from being counted twice.


The complete system contains forms of several quite different geometric meanings. Some express
special contacts of the quartic with lines, others express the presence of singular polar curves,
and still others measure degenerations of the net of conics associated with the quartic. The
invariant-theoretic calculation itself does not require these interpretations. Nevertheless they
provide a useful check: when the symbolic form is specialized to a quartic with a known singularity,
the corresponding invariant must vanish with the predicted multiplicity.


The final theorem can be stated briefly. The concomitants displayed in the tables generate the full
algebra of concomitants of the ternary quartic. Every concomitant of the ground form is a polynomial
expression in them. Moreover the stated symbolic relations suffice to reduce any symbolic expression
to such a polynomial. The proof is by induction on coefficient degree, using the reduction lemmas
and the dimension comparison at each bidegree. This completes the construction of the form system.


\subsection{Additional reductions and formulation}

The enumeration is best read as an induction on the coefficient degree. In degree one there is only
the ground form and its polar relatives. In degree two the determinant factors first appear, and the
possible concomitants are still few enough to list directly. In higher degrees the number of
symbolic monomials is large, but all monomials have a finite signature: the number of determinant
brackets, the remaining powers of each symbolic letter, and the distribution of point and line
variables. The reduction preserves this signature except when it replaces a product by forms of
lower degree. This is why the induction terminates.

At several places the calculation uses what may be called the exchange of letters. A symbolic
product with letters \(a,b,c,d\) may be expanded by a determinant identity and rewritten as a sum in
which one of the letters occurs in a simpler position. Repeating the exchange eventually brings
every product to a normal arrangement. The signs introduced by the exchanges are fixed by the
orientation of the determinants. Once a normal arrangement is chosen, two symbolic products can be
compared unambiguously.

The Cayley series supplies an independent verification of the enumeration. For each type \((d,r,s)\)
it gives the number of independent concomitants that can exist. The table of products generated from
earlier forms gives a lower bound on that number, and the symbolic reductions give an upper bound.
When the two bounds agree, the entry is closed. When they do not agree, the difference is exactly
the number of new primitive forms to be inserted. This is the numerical mechanism behind the phrase
complete system.

A point that is easy to miss in the printed tables is that the same concomitant can occur in several
symbolic disguises. For instance, a product obtained by first polarizing and then contracting can
coincide with one obtained by first multiplying and then applying a transvection. The dissertation
resolves this by retaining one representative and recording the identity that identifies the others
with it. Thus the table is not only a list of generators; it is also a record of the reductions
among candidate generators.

The specialization argument is also essential. A symbolic identity proved for a general quartic is
an identity in the coefficients and therefore remains true for a singular quartic. Conversely, the
independence count is made in the generic case; if a generator specializes to zero for a special
quartic, this does not remove it from the complete system. The system is an algebraic system over
the coefficient ring, not a minimal system for each individual member of the family.

The complete system therefore has two roles. It gives a finite algebra of concomitants for
calculations with a fixed ternary quartic, and it gives a coordinate algebra for the quotient
problem in which quartics are considered up to projective transformation. The second role explains
why it is not enough to know invariants alone. Covariants and mixed concomitants encode the geometry
of associated points, lines, and polar curves; they are needed to reconstruct the full projective
information.


\subsection{Further proof details}

A final check concerns relations among the fundamental forms themselves. The theorem that a list
generates all concomitants is distinct from the problem of finding all syzygies among the
generators. The dissertation uses only those syzygies needed to prove generation. Some relations are
displayed because they are required for reductions; others are consequences of the symbolic
determinant identities but are not needed for the proof of completeness. Thus the printed system is
a generating system, not a free system.

The low-degree entries also serve as anchors for the induction. If an alleged reduction in higher
degree uses a product of lower forms, each of those lower forms must already have been normalized.
Otherwise the same expression could be counted under two different names. This is why the tables
proceed in a fixed order and why every operation is recorded at the point where it first becomes
available.

Geometrically, the ternary quartic is a plane curve of degree four. Its invariants detect projective
features of the curve; its covariants produce associated curves and point configurations. The
symbolic system therefore encodes both the algebra of the coefficients and the projective geometry
of the curve. The computation is long because these two aspects are interwoven, not because the
principles are obscure.


\section{5. Rational function fields}

\textit{Original publication: Jahresbericht der Deutschen Mathematiker-Vereinigung 22 (1913), pp. 316--319.}

Let \(K\) be a rational function field over a ground field \(k\). The question considered here is
how the property of being rational can be recognized without reference to a special system of
independent variables. If \[   K=k(x_1,\ldots,x_n), \] then every other system of algebraically
independent generators gives the same field. The problem is to determine which subfields and
extensions obtained from such a field are again rational and how this can be detected by intrinsic
algebraic operations.


A first observation is elementary but useful. If \(u_1,\ldots,u_n\) are algebraically independent
and \(K\) is algebraic over \(k(u_1,\ldots,u_n)\), then equality need not hold. Rationality is
therefore stronger than having the correct transcendence degree. One must ask for a transformation
whose inverse is again rational. Thus the central object is the group of birational substitutions of
the variables, or equivalently the automorphism group of \(k(x_1,\ldots,x_n)\) over \(k\).


In one variable the theorem is classical: every subfield of \(k(x)\) of transcendence degree one is
again a rational function field. This is the algebraic form of the theorem that the projective line
has genus zero. In several variables the analogous statement is false in this simple form.
Nevertheless many subfields arising from linear substitutions or from finite groups of substitutions
remain rational, and the proof can often be reduced to a choice of suitable invariants.


Suppose first that a finite group \(G\) acts linearly on \(k(x_1,\ldots,x_n)\), permuting or
linearly transforming the variables. The fixed field \[   K^G=\{z\in K: z^S=z\text{ for all }S\in
G\} \] has the same importance for rational function fields that invariant rings have for polynomial
forms. The rationality question asks whether \(K^G\) can be written as \(k(y_1, \ldots,y_n)\). In
favorable cases the elementary symmetric functions, or their analogues for the given representation,
give such a system.


The method used here is to separate the algebraic and the transcendental parts of the argument.
Choose a separating transcendence basis of the fixed field. Then the original field is algebraic
over this fixed field; its degree is bounded by the order of the group. If one can show that the
chosen invariants already give an extension of degree equal to the group order, then they form a
full rational system. Thus a field-theoretic degree calculation replaces the direct construction of
an inverse substitution.


For permutation groups this gives the familiar theorem immediately. If \(G\) is the full symmetric
group on the variables, the elementary symmetric functions generate the fixed field. For subgroups
the situation is subtler. One may need resolvent functions whose conjugates under the symmetric
group separate the cosets of \(G\). The corresponding fixed field is then generated by the symmetric
functions together with one resolvent root. Rationality is obtained only when this root can be
replaced by a rational parameter system.


The note emphasizes that rationality is invariant under extension of the ground field only in one
direction. If a field becomes rational after adjoining constants, it does not necessarily follow
that it was rational over the original field. The obstruction is a descent obstruction. In the
language used later, it is expressed by the failure of a system of conjugate generators to be
definable over the smaller field. Even when the algebraic relations have coefficients in \(k\), the
generators themselves may be defined only over an extension.


This distinction is important for invariant theory. A complete system over an algebraically closed
field may descend only after replacing its generators by invariant combinations. Conversely,
rational generators over a splitting field may satisfy conjugacy relations that prevent them from
giving rational generators over \(k\). Thus the rationality of the field of invariants is not a
purely geometric question; it also involves arithmetic of the constants.


The guiding lemma is the following. Let \(L/K\) be a finite Galois extension and let \(x\) be
transcendental over \(L\). If the group acts on \(L(x)\) by fractional linear transformations in
\(x\), then the fixed field is rational over the fixed field of \(L\), provided that a fixed divisor
of degree one exists. In elementary terms, one constructs a parameter by taking a quotient of two
semi-invariant linear forms. The proof is just the transformation law for binary linear forms,
rewritten in field language.


The conclusion is not a classification theorem, but a collection of reductions. To decide whether a
field is rational one should first isolate the finite algebraic part, then test whether the
remaining transcendental variables can be chosen invariantly. This method will be used again in the
more systematic treatment of fields and systems of rational functions. The present note records the
field-theoretic core of that treatment.


\subsection{Additional reductions and formulation}

The one-variable theorem may be used as a model for the general method. If \(F\) is a subfield of
\(k(x)\) of transcendence degree one, choose an element \(u\in F\) of minimal degree as a rational
function of \(x\). Any other element of \(F\) is algebraic over \(k(u)\), but the minimality forces
the degree of \(k(x)\) over \(k(u)\) to coincide with the degree over \(F\). Hence \(F=k(u)\). In
higher dimension no such simple minimal degree argument exists, and the failure of this argument is
precisely where the new difficulties begin.

For finite groups the fixed field can be approached through the normal basis idea. A generic element
\(z\) has conjugates \(z^S\) under the group. The elementary symmetric functions of these conjugates
lie in the fixed field, and \(z\) satisfies the corresponding equation. If the conjugates are all
distinct, the degree of the equation is the group order. The fixed field is then obtained by
adjoining to the symmetric fixed field enough coefficients to describe the orbit. This is the field
version of the resolvent construction.

The rationality of a fixed field is often proved by finding triangular generators. One first chooses
an invariant involving all variables, then replaces one variable by a quotient invariant, and
proceeds by induction. In permutation examples this produces elementary symmetric functions. In
affine examples it produces differences and sums. The important point is that each step must be
reversible by rational operations; otherwise one has only algebraic generation, not rational
generation.

Descent over the ground field can be stated concretely. Suppose \(L/k\) is a finite extension and
\(L\otimes_k F\) is rational over \(L\). Choose rational generators \(y_i\). For each automorphism
\(S\) of \(L/k\), the elements \(S(y_i)\) form another rational generating system. Descent asks
whether there is a system of rational functions in the \(y_i\), invariant under all \(S\), that
still generates the same field. This is generally a non-trivial problem of simultaneous rational
transformation.

Thus even the modest question whether a field is rational leads to three layers: algebraic
dependence, birational inversion, and descent of constants. The first is decided by prime ideals of
relations. The second is decided by the degree of the field over a chosen transcendence basis. The
third is decided by the Galois action on possible bases. Keeping these layers separate prevents
several common mistakes in the use of rational parameters.


\subsection{Further proof details}

One practical test for rationality is the construction of a Lüroth-type parameter on a divisor or
pencil. If a family of hypersurfaces has a rational member through the general point and the
residual intersection is controlled, the quotient of two equations of the pencil may serve as a
parameter. This geometric argument translates into the field statement that a subfield generated by
a pencil has transcendence degree one and is therefore rational.

In several variables, the failure of rationality is often invisible if one looks only at general
points. A unirational field can be covered by rational parameters without being rational itself.
Field-theoretically this means that the ambient rational field is finite over the subfield, but the
degree is greater than one. The rational functions give a dominant parametrization; they do not give
a birational inverse. This distinction is one of the motivations for the later systematic use of
field degrees.

The note's conclusions are deliberately modest. It does not claim that every fixed field of a finite
group is rational. Rather, it isolates the constructions that do give rational fixed fields and
indicates where the obstruction lies. This formulation later became one of the standard forms of
Noether's problem: decide whether the field of invariants of a rational field under a finite group
is again rational over the ground field.


\section{6. Fields and systems of rational functions}

\textit{Original publication: Mathematische Annalen 76 (1915), pp. 161--196.}

Let \(K\) be a field of rational functions in finitely many variables over a ground field \(k\). A
system of rational functions \[   f_1,\ldots,f_m\in K \] defines a subfield \(k(f_1, \ldots,f_m)\).
The basic problem is to determine when two systems define the same subfield, when one system can be
transformed into another by rational substitutions, and when a subfield so obtained is itself purely
transcendental. These questions are the natural field-theoretic counterpart of elimination theory.


The first step is to replace a system of functions by the ideal of algebraic relations among them.
If \(Y_1, \ldots,Y_m\) are indeterminates, the kernel of the homomorphism \[
k[Y_1,\ldots,Y_m]\longrightarrow K,   \quad Y_i\mapsto f_i, \] is a prime ideal. Its quotient field
is the field \(k(f_1, \ldots,f_m)\). Thus the study of systems of rational functions is equivalent
to the study of prime ideals together with their quotient fields. This formulation makes possible
the use of dimension and specialization arguments.


A subsystem \(f_{i_1}, \ldots,f_{i_r}\) is called independent if there is no algebraic relation
among its members. A maximal independent subsystem has cardinality equal to the transcendence degree
of the generated field. The remaining functions are algebraic over the field generated by such a
subsystem. However, different maximal independent subsystems need not give isomorphic algebraic
extensions. The dependence on the chosen basis is controlled by the field degree of the total system
over that basis.


The main technical device is normalization by a linear change of variables. After replacing the
functions by general linear combinations, one may assume that the first \(r\) functions form a
separating system and that the remaining functions are integral over the polynomial ring generated
by them, after clearing denominators. This is a rational-function version of Noether normalization.
It permits all later arguments to be carried out over a polynomial subring, where denominators and
exceptional specializations can be controlled.


Let \(F=k(f_1, \ldots,f_m)\). A rational representation of \(F\) is a choice of algebraically
independent elements \(u_1, \ldots,u_r\) such that \(F\) is algebraic over \(k(u_1, \ldots,u_r)\).
The degree of this algebraic extension is the order of the representation. If the order is one, the
representation is parametric: the field is rational. The paper shows how to test the order by
eliminating the original variables and computing the degree of the resulting prime ideal over the
chosen parameters.


The role of resultants is purely auxiliary. A resultant may eliminate variables and produce an
equation for one function over the others, but the resultant often contains extraneous factors
introduced by denominators or by components at infinity. The correct object is the prime ideal of
relations. Resultants must be saturated with respect to the product of the denominators and then
decomposed. Only the component corresponding to the given field is relevant.


The same point is visible in geometric language. The rational functions define a rational map from
affine or projective space into another space. The field \(F\) is the function field of the image.
Exceptional divisors and base loci may alter equations in the image if they are not removed. Working
with prime ideals and quotient fields suppresses these accidental components. Birational equivalence
is then equality of function fields, not equality of one particular set of equations.


A useful criterion is obtained from specialization. Suppose the coefficients of the functions depend
on parameters. If for one specialization the functions are algebraically independent and the
relevant degree is \(d\), then outside a proper exceptional set the same independence and the same
degree hold. This is because the leading coefficients and discriminants that control the integral
equations vanish only on a proper algebraic subset. Thus generic rationality questions can be
reduced to one well-chosen specialization.


The treatment also separates absolute and relative rationality. A field may be rational over an
algebraic extension of \(k\) but not over \(k\) itself. If rational generators over the larger field
are permuted by the Galois group, the descent question becomes the problem of finding invariant
combinations that remain algebraically independent and generate the same field. This is the same
obstruction encountered in the theory of rational function fields and will later reappear in
invariant theory and Galois theory.


The principal theorem can be formulated as follows. Every field generated by a finite system of
rational functions admits a finite normal representation over a purely transcendental subfield;
after a general linear transformation the generators may be chosen among linear combinations of the
original functions. Equality of two such fields is decided by equality of the corresponding prime
ideals after passage to quotient fields. Rationality is the special case in which the normal
representation has degree one.


In applications, one begins with explicit functions and forms the relation ideal. One then chooses a
maximal independent subsystem and computes the degree of the remaining algebraic extension. If the
degree is one, the system is rational and the chosen subsystem is a rational parameter system. If
the degree is larger, no rational inverse exists for that particular subsystem; one may still try
another subsystem or a birational transformation, but the obstruction is now a genuine field degree.


The significance of the paper lies in the replacement of formula manipulation by field-theoretic
invariants. A rational system is not an array of formulae but a finitely generated field extension.
Its essential data are transcendence degree, algebraic degree over a normalizing subfield, and the
prime ideal of relations. These data survive changes of variables and make the theory independent of
accidental representations.


\subsection{Additional reductions and formulation}

The proof of normalization is elementary in spirit. Let the functions be written with common
denominator, and let the numerator variables be subjected to a general linear transformation. The
bad transformations are those for which some leading coefficient vanishes or some algebraic
dependence degenerates. These conditions are algebraic equations in the transformation coefficients.
Since they do not hold identically, a transformation outside their zero set has the required
property. This is the algebraic content of the word general.

When the relation ideal is prime, the quotient ring still may not be integrally closed.
Normalization of the quotient ring corresponds to passing from a singular image variety to its
normal model. At the level of function fields this passage changes no rational functions, but it may
simplify integral equations. Noether uses this distinction systematically: equations describe a
model, whereas the quotient field describes the birational object itself.

Equality of two function fields generated by two systems \(f_i\) and \(g_j\) can be checked by
mutual rational expression. It is enough to show that every \(f_i\) is rational in the \(g_j\) and
every \(g_j\) is rational in the \(f_i\). Algebraically this means that the two embeddings of
polynomial rings into the ambient field have quotient fields with the same image. The condition is
unaffected by clearing denominators, provided that the denominators are not allowed to vanish
identically.

The degree of a rational map appears here as a field degree. If the image field \(F\) is generated
by \(r\) independent functions and \([K:F]=d\), then a general point of the image has \(d\)
preimages in the original parameter space, counted algebraically. Birationality is exactly the case
\(d=1\). This gives a purely algebraic criterion for when a system of rational functions can serve
as a change of variables.

The paper's method is especially useful when parameters occur. A family of systems defines a prime
ideal over a parameter ring after removing exceptional components. The degree over a normalizing
subfield is constant on a dense open set of the parameter space. Hence properties such as
independence, separability, and birationality are generic. Special parameter values may be
exceptional, but they do not alter the general theorem.

In later ideal-theoretic language one would say that the coordinate ring of the image is a finitely
generated domain and that the original rational field is a finite extension of its field of
fractions after normalization. The 1915 treatment does not yet use the later vocabulary of primary
decomposition in its final form, but it already uses the same strategy: replace computations with
equations by computations with prime ideals and quotient fields.


\subsection{Further proof details}

The relation ideal also records implicit equations. Given a rational parametrization, eliminating
the parameters gives equations for the image. But the implicit equations alone do not determine
whether the parametrization is one-to-one. The field degree over the image field measures this. Thus
implicitization and inversion are separate problems: the former belongs to ideals, the latter to
quotient fields.

A useful way to organize the calculations is to choose a tower \[   k(u_1, \ldots,u_r)\subset
F\subset K. \] The first extension is finite by normalization, and the second is the inclusion of
the generated field into the ambient rational field. If \(F=K\), the system gives a birational
coordinate system. If \(F\neq K\), the missing functions represent the algebraic ambiguity of the
parametrization. The tower makes the ambiguity explicit.

The paper also anticipates the language of generic freeness. After excluding a finite set of bad
denominators, modules that occur in the calculation become free over the parameter ring, and ranks
become constant. This justifies specialization arguments: a rank or degree computed at one good
specialization is the generic rank or degree. The exceptional cases are not ignored; they are
precisely the zeros of the denominators and discriminants introduced during normalization.


\section{9. The most general domains of completely transcendental numbers}

\textit{Original publication: 1916.}

Consider a domain generated from a ground field by elements subject to no algebraic relations except
those forced by the construction. Such elements are called completely transcendental. The question
is to describe the most general domains in which these elements can be manipulated as independent
quantities while still allowing algebraic adjunctions, specialization, and passage to quotient
fields. This is a foundational problem for the algebraic use of parameters.


Let \(k\) be a base field and let \(x_1, \ldots,x_n\) be algebraically independent. The polynomial
ring \(k[x_1, \ldots,x_n]\) is the simplest domain of completely transcendental quantities, and its
quotient field is \(k(x_1, \ldots,x_n)\). But many constructions require adjoining further elements
integral or algebraic over this ring, or localizing at systems of denominators. The problem is to
determine which properties of independence survive these operations.


A domain \(R\) over \(k\) is to be viewed through its finite subextensions. Every element of \(R\)
should lie in a subdomain generated by finitely many transcendental elements and finitely many
algebraic elements over them. This finiteness condition prevents pathologies and permits the use of
elimination. Thus even a very large domain is controlled by its finitely generated subdomains.


The essential invariant is the number of independent transcendental generators. If \(R\) has
quotient field \(K\), this number is the transcendence degree of \(K/k\). It is independent of the
chosen generators. If \(u_1, \ldots,u_r\) and \(v_1, \ldots,v_s\) are two transcendence bases of
\(K/k\), then \(r=s\). Moreover each set can be transformed into the other through a finite chain of
elementary replacements, at least after restricting to suitable finitely generated subfields. This
gives a concrete algebraic meaning to the dimension of the domain.


Specialization is handled by homomorphisms from the polynomial subdomain to algebraic extensions of
\(k\). An algebraic relation that holds after all admissible specializations must already hold in
the original domain. Conversely, a denominator that is non-zero in the general domain may vanish
under exceptional specializations. Therefore one must always exclude a finite union of hypersurfaces
when using specialization to prove general statements.


The paper insists on the difference between formal independence and algebraic independence in a
fixed field. Formal variables are symbols; algebraically independent elements are elements of an
extension field satisfying no polynomial relation. Any identity in formal variables remains true
after substitution by algebraically independent elements. But a statement involving denominators or
algebraic adjunctions may require restricting the substitutions. This is the reason for introducing
domains rather than only fields.


The construction of the most general domain proceeds by taking a universal polynomial domain and
then factoring only by the prescribed algebraic relations. If no relations are prescribed, the
domain remains a polynomial ring. If algebraic elements are prescribed by equations, one forms the
corresponding quotient and then selects the component determined by the desired prime ideal. The
quotient field of that component is the most general field satisfying the requirements.


This language clarifies the use of generic points. A generic point of an irreducible variety is not
a mysterious geometric object; algebraically it is the homomorphism from the coordinate ring into
its quotient field. The coordinates of the generic point are completely transcendental modulo the
prime relations defining the variety. Statements about the general point are statements in this
quotient field. Exceptional cases correspond to specializations where denominators vanish or
algebraic degrees drop.


The final conclusion is that domains of completely transcendental numbers are best treated as
universal domains built from polynomial rings, prime ideals, localizations, and finite algebraic
extensions. Their generality is not obtained by allowing infinitely many uncontrolled elements, but
by requiring that every finite calculation take place in a finitely generated subdomain with a
specified transcendence basis. This is the algebraic setting in which later normalization and
ideal-theoretic arguments become possible.


\subsection{Additional reductions and formulation}

A recurring question is whether a construction depends on the names of the transcendental elements.
It should not. If \(x_1,\ldots,x_n\) and \(u_1,\ldots,u_n\) are two transcendence bases of the same
rational field, any identity that is invariant under field isomorphism must have the same meaning in
both systems. The formalism of universal domains is designed to express statements in this
coordinate-free way, while retaining enough coordinates for actual calculation.

The passage from polynomial rings to quotient fields must be made with care. In the polynomial ring,
divisibility and specialization are visible. In the quotient field, every non-zero polynomial has
become invertible, and exceptional cases have disappeared. A proof that uses quotient fields may
therefore prove only a generic statement. To recover a statement about specializations one must
return to a localized polynomial ring and record which denominators have been inverted.

Algebraic adjunctions are handled by prime ideals in polynomial rings over the transcendental base.
If an element \(y\) is to satisfy an irreducible equation \(F(y)=0\), the universal domain is the
quotient by the prime ideal generated by the chosen irreducible factor. If the equation factors
after extending constants, the different factors correspond to different branches. This is the
algebraic origin of the distinction between absolute and relative irreducibility.

The independence theorem for transcendence bases underlies the entire construction. If \(K\) is
finitely generated over \(k\), every maximal algebraically independent subset has the same
cardinality. Furthermore \(K\) is algebraic over the rational field generated by such a subset. Thus
the entire domain can be built in two stages: a completely transcendental stage and an algebraic
stage. This two-stage decomposition is the model for many later normalization arguments.

The most general domain for a problem is therefore not unique as a ring with named generators, but
it is unique up to the evident universal property. Any other domain satisfying the same formal
relations receives a specialization from it, after excluding denominators. This universal property
is what justifies calculations with generic coefficients and generic roots. The generic calculation
is a calculation in an initial algebraic situation from which special cases are obtained by
homomorphism.


\subsection{Further proof details}

The terminology completely transcendental is meant to exclude hidden algebraic relations. If a
calculation with parameters produces a polynomial relation among them, then the parameters were not
completely transcendental for that problem. The universal domain is constructed so that every
relation is visible in the defining prime ideal. This is the safest setting for proofs by generic
specialization.

One can also use these domains to compare different problems. If two universal domains have
isomorphic quotient fields over the base field, then the two generic situations are birationally
equivalent. If the isomorphism preserves the distinguished subrings, the equivalence is stronger and
respects specialization. Thus the language of domains distinguishes birational equivalence from
isomorphism of chosen models.

The later ideal theory will replace many of these arguments by statements about noetherian rings and
integral closure. Still, the guiding idea remains the same: every general algebraic situation should
be represented by a finitely generated domain whose quotient field is the field of the generic case.
Once this is done, questions about specialization, dependence, and algebraic extension can be
treated uniformly.


\section{10. Functional equations of the isomorphic mapping}

\textit{Original publication: 1916.}

Let two algebraic domains be given by generators and relations. An isomorphic mapping between them
is determined by the images of the generators, but those images must satisfy the same relations.
Thus the problem of isomorphism can be written as a system of functional equations. The aim is to
express this system algebraically and to determine when its solutions form a finite or rational
family.


Let \[   R=k[x_1, \ldots,x_n]/\mathfrak p,   \qquad   R'=k[y_1, \ldots,y_m]/\mathfrak q \] be
integral domains. A homomorphism \(R\to R'\) is defined by choosing elements \(X_i(y)\in R'\) such
that every relation in \(\mathfrak p\) vanishes after substitution \(x_i\mapsto X_i(y)\). If the
homomorphism is to be an isomorphism, there must also be inverse expressions \(Y_j(x)\) satisfying
the corresponding relations and the two systems of composition equations.


The functional equations therefore consist of three parts: preservation of the defining relations of
\(R\), preservation of the defining relations of \(R'\), and inversion equations. Written with
undetermined coefficients for the functions \(X_i\) and \(Y_j\), these become polynomial equations
in those coefficients. Thus the set of isomorphisms is itself represented by an algebraic system
whenever degree bounds are imposed.


Degree bounds are not arbitrary. In many problems they follow from dimension, from the degree of the
defining equations, or from the requirement that certain divisor classes be preserved. Once a bound
is known, the existence of an isomorphism is an elimination problem. One eliminates the coefficients
of the inverse functions and obtains conditions on the original domains. Conversely, a solution of
the coefficient equations gives an actual isomorphism after checking the denominators.


The method applies particularly well to rational function fields. An automorphism of \(k(x)\) is a
fractional linear transformation, and the functional equations reduce to the invertibility of a
two-by-two determinant. In several variables the Cremona transformations appear. They are again
described by polynomial equations among the coefficients of rational functions and their inverses.
The complexity lies not in the concept of isomorphism but in obtaining useful degree bounds.


If a group is prescribed to act by isomorphisms, the functional equations acquire additional
composition relations. For each group element \(S\) one introduces a system of functions \(X_i^S\).
The equations express \[   X^{ST}=X^S\circ X^T,   \qquad X^1=X. \] Together with the defining
equations of the domain, these are the algebraic equations for a representation of the group by
automorphisms. This viewpoint anticipates the treatment of equations with prescribed group.


The paper also makes clear why isomorphism is a relative notion. If the ground field is enlarged,
new solutions of the functional equations may appear. An isomorphism over \(\bar k\) need not
descend to \(k\). The descent condition is the invariance of the solution under the Galois group,
possibly after conjugation by automorphisms of the object. Thus the functional equations describe
both the geometric isomorphisms and the arithmetic obstruction to their definition over the ground
field.


In modern notation one would say that the isomorphisms form a functor. The paper treats the same
idea by direct equations. For every coefficient domain \(A\) one may ask for solutions with
coefficients in \(A\); substitution and composition are polynomial operations, hence compatible with
homomorphisms of \(A\). The representing equations therefore encode the whole family of
isomorphisms, not only the individual transformations over an algebraically closed field.


The conclusion is that the problem of recognizing isomorphic algebraic structures can be brought
under elimination theory once the transformations are bounded. The functional equations are
explicit, functorial, and stable under specialization. In this form the theory provides a bridge
between rational substitutions, automorphism groups, and the Galois-theoretic construction of
equations.


\subsection{Additional reductions and formulation}

One should distinguish an isomorphism from a correspondence that is generically one-to-one. A
rational map may have an inverse on dense open subsets but fail to be a regular isomorphism on the
chosen affine models. The functional equations detect the birational inverse, while additional
denominator conditions detect where the inverse is defined. This distinction is why the equations
are naturally written in function fields first and only afterward translated into coordinate rings.

If the domains carry extra structure, such as a distinguished subfield, a grading, or a divisor at
infinity, the functional equations acquire preservation conditions. For example, a graded
isomorphism must send homogeneous elements of degree one to homogeneous elements of degree one. Such
restrictions often provide the degree bounds that are otherwise missing. The general method is
therefore flexible: one writes equations not only for multiplication but also for every structure
required to be preserved.

Composition of two candidate isomorphisms is polynomial in the coefficients after clearing
denominators. The equations expressing that the composition is the identity are therefore algebraic
equations. This is an important finiteness point: invertibility is not an infinite condition once a
degree bound has been imposed. It is represented by finitely many coefficient equations obtained by
comparing the numerator and denominator of each composed coordinate.

For automorphism groups, the same equations define a group law on the solution set. The identity
solution and inverse solution are given by polynomial operations in the coefficients, again after
localization by the determinant or resultant that expresses non-degeneracy. Thus the automorphisms
form an algebraic group or an algebraic family in the cases where the degree is bounded. The later
language of group schemes is not used, but the construction is already present in equation form.

The descent problem can be solved in favorable cases by averaging only after the isomorphism
equations have been linearized. Direct averaging of rational transformations usually fails because
denominators do not average well. Instead one studies the induced action on a finite-dimensional
vector space of forms of bounded degree. If an invariant point satisfying the non-degeneracy
condition exists in that space, it gives a descended isomorphism. This is the algebraic content
behind many classical descent arguments.


\subsection{Further proof details}

The same functional-equation method applies to equivalence rather than isomorphism. If objects are
considered equivalent after multiplying equations by units, changing coordinates, or adding
redundant generators, the equations are enlarged by the corresponding auxiliary unknowns.
Equivalence is then expressed by the solvability of this larger system. This is how classification
problems enter the same algebraic framework.

A point of the construction is that automorphisms are not first found one by one. They are found as
the solution set of a finite algebraic system. This makes it possible to speak of specializations in
which the automorphism group jumps. The generic automorphism group may be smaller than that of a
special member, and the jump is detected by additional solutions of the functional equations after
specialization.

The inversion equations are usually the most restrictive part. A homomorphism of coordinate rings is
easy to write; proving that it has an inverse forces relations among the coefficients of the
proposed images. In rational form these relations include non-vanishing of determinants or
resultants. The algebraic equations describe the closed part of the condition; the non-vanishing
conditions describe the open part.


\section{11. Equations with prescribed group}

\textit{Original publication: 1918.}

Let \(G\) be a finite group. The problem is to construct equations whose group is \(G\), and to do
so in a form that separates the formal group action from accidental specializations. This is the
inverse problem in its algebraic form. One begins with a rational function field on which \(G\) acts
faithfully, forms the fixed field, and then regards the original field as a Galois extension of that
fixed field.


Let \(x_g\) be independent variables indexed by the elements of \(G\), and let \(G\) act by left
translation, \[   S(x_g)=x_{Sg}. \] The field \(K=k(x_g:g\in G)\) is then a regular representation
field for \(G\). Its fixed field \(K^G\) has the property that \(K/K^G\) is a Galois extension with
group \(G\). A primitive element \(z\in K\) gives an equation over \(K^G\) whose roots are the
conjugates \(S(z)\). This equation has group \(G\) in the generic sense.


The difficulty is to replace the many variables by as few parameters as possible and to decide when
specialization preserves the group. The coefficients of the generic equation lie in the fixed field.
If this fixed field is rational, the equation can be written with independent parameters. If it is
not known to be rational, one still obtains an equation over the fixed field, but not necessarily
over a purely transcendental parameter field. Thus the rationality problem for fixed fields enters
directly into the inverse Galois problem.


Specialization is controlled by the discriminant and by the resolvents corresponding to proper
subgroups. If the discriminant is non-zero, the specialized equation is separable. If no resolvent
belonging to a proper subgroup acquires a rational root, the group does not drop into that subgroup.
Since there are only finitely many proper subgroups, the exceptional specializations are contained
in a finite union of proper algebraic sets. Hence the general specialization retains the prescribed
group.


The construction is invariant under change of primitive element. If \(z\) and \(z'\) both generate
\(K\) over \(K^G\), their equations have the same splitting field and the same group. The
coefficients are different functions on the same parameter domain. Thus the group belongs to the
extension, not to the particular polynomial. This distinction is essential when comparing equations
obtained by different resolvent constructions.


The regular representation may be replaced by any faithful representation of \(G\) by substitutions
of rational functions. If the action on a rational field \(k(t_1, \ldots,t_n)\) is faithful, then
the extension over the fixed field again has group \(G\). A smaller representation can yield
equations with fewer parameters. But the fixed field may become harder to identify. The trade-off is
between the number of parameters and the rationality of the invariant field.


If \(G\) is a subgroup of the symmetric group \(S_n\), one can take the roots themselves as
variables. The elementary symmetric functions give the full symmetric fixed field, and additional
resolvents distinguish \(G\) from its overgroups. This is the classical approach. The
field-theoretic formulation shows that it is a special case of the same invariant construction, with
the roots providing a permutation representation.


The theorem may be summarized as follows. For every finite group \(G\) there is a generic Galois
extension with group \(G\), obtained from a rational function field carrying the regular action of
\(G\). Its specializations, outside a proper exceptional set and subject to the usual separability
condition, give equations whose group is contained in \(G\); outside the additional resolvent loci
the group is exactly \(G\). The remaining problem is to make the parameter field rational and
minimal.


This conclusion shifts the inverse problem from the existence of a formal equation to the
rationality and descent of invariant fields. It explains why the problem is easy to formulate for
every finite group but difficult to solve over a prescribed ground field with few parameters. The
group action is elementary; the arithmetic of the fixed field is the subtle part.


\subsection{Additional reductions and formulation}

The generic equation obtained from the regular representation is universal in a weak but useful
sense. Any Galois extension with group \(G\) can be obtained by specializing the independent
variables to a normal basis, at least after adjoining the necessary constants and avoiding the
discriminant. The coefficients of the generic equation then specialize to the coefficients of the
given equation. This explains why the regular representation is the natural starting point despite
its many variables.

Resolvents measure subgroup containment. If \(H\subset G\), choose a rational function fixed by
\(H\) and by no larger subgroup. Its conjugates under \(G\) give a resolvent equation whose rational
root after specialization is equivalent to the Galois group lying in a conjugate of \(H\). Thus the
group is exactly \(G\) if none of the resolvents for maximal proper subgroups has a rational root.
This is the classical resolvent criterion in invariant-field language.

The construction also shows why different polynomials may have the same group but different
parameter fields. Changing the primitive element changes the embedding of the extension into a
polynomial equation. The splitting field and its group remain the same, but the coefficients may
generate a smaller or larger subfield of the invariant field. The best equation is often the one
whose coefficients generate the smallest rational subfield over which the extension is still
regular.

When \(G\) has a faithful linear representation of small dimension, one may replace the regular
variables by coordinates in that representation. If the invariant field is rational, this gives a
generic equation with fewer parameters. The obstruction is exactly the rationality problem for
\(k(V)^G\). Thus Noether's problem on rationality of invariant fields appears not as a side issue
but as the central compression problem for equations with prescribed group.

The exceptional specialization set is not merely the discriminant. The discriminant prevents
collisions among roots; resolvent conditions prevent the group from falling. There can also be
denominator conditions coming from the chosen primitive element and from the chosen rational
generators of the fixed field. The correct statement is therefore always: outside a finite union of
proper algebraic exceptional sets, the specialization has the prescribed group.


\subsection{Further proof details}

For a subgroup chain \(H\subset G\), resolvents can be arranged compatibly. A root of the
\(H\)-resolvent gives a reduction of the group to a conjugate of \(H\); further resolvents then test
smaller subgroups. In this way the lattice of subgroups is reflected in a tower of resolvent
equations. The generic construction contains all these resolvents at once, because every
intermediate fixed field of \(K/K^G\) corresponds to a subgroup of \(G\).

The field of constants remains a serious issue. Over an algebraically closed field the regular
representation construction is formal. Over a smaller field, the desired group action and the
primitive element must be defined over that field. If the group representation requires roots of
unity or other constants, the resulting equation over the smaller field may have a different
arithmetic behavior. This is another form of the descent problem.

The final shape of the theorem is therefore twofold. Universally, every finite group occurs as the
group of a regular extension of a suitable invariant field. Arithmetically, the production of
equations over a specified rational parameter field depends on rationality and descent. The
distinction prevents the formal construction from being mistaken for a complete solution of the
inverse Galois problem over a given base.


\section{12. Invariants of arbitrary differential expressions}

\textit{Original publication: Nachrichten der Gesellschaft der Wissenschaften zu Goettingen, 1918.}

Let a differential expression be formed from dependent variables and their derivatives up to some
finite order. A group of transformations of the independent and dependent variables acts not only on
the variables themselves but on all derivatives, by the prolonged transformation. An invariant
differential expression is one that remains unchanged under this prolonged action. The purpose is to
describe such invariants without restricting to special forms of the expression.


If \(x=(x_1, \ldots,x_n)\) are independent variables and \(u=(u_1, \ldots,u_m)\) dependent
variables, then derivatives are denoted by \(u^\alpha_j\), where \(\alpha\) is a multi-index. A
finite-order differential expression is a function of finitely many such jet variables. A
transformation group acts on the jet variables by the chain rule. The invariance condition is
therefore a system of differential equations for the expression.


For an infinitesimal group with generators \[   X_\nu=\sum_i \xi_{\nu i}\frac{\partial}{\partial
x_i}         +\sum_j \eta_{\nu j}\frac{\partial}{\partial u_j}, \] the prolonged generators
\(X_\nu^{(r)}\) act on derivatives up to order \(r\). A differential expression \(I\) of order \(r\)
is invariant precisely when \[   X_\nu^{(r)} I=0 \] for all generators. Thus the determination of
invariants becomes the integration of a complete system of first-order linear partial differential
equations on jet space.


The completeness of this system follows from the closure relations of the infinitesimal group. If
\([X_\mu,X_\nu]\) is a linear combination of the generators, then the same is true after
prolongation. Consequently the equations \(X_\nu^{(r)}I=0\) are compatible. Their solutions are
functions of a maximal set of independent differential invariants. The number of such invariants is
the dimension of the jet space minus the rank of the prolonged group action.


Invariant differentiation is the operation that carries invariants to new invariants. An ordinary
derivative of an invariant need not be invariant, because the independent variables themselves
transform. One must form special differential operators \[   D_1, \ldots,D_s \] which commute with
the prolonged group action up to linear combinations. Applying these operators to known invariants
yields further invariants. This is the differential analogue of forming covariants from algebraic
invariants by transvection.


The finiteness question is more delicate than in polynomial invariant theory, since differentiation
produces infinitely many expressions. The appropriate statement is not that there are only finitely
many invariants, but that all invariants can be generated from finitely many fundamental
differential invariants by invariant differentiation and algebraic operations. The proof depends on
the stabilization of the ranks of the prolonged actions and on the ability to solve for higher
derivatives in terms of invariant derivatives.


The treatment is independent of the special nature of the differential expression. It applies to
scalar functions, systems of functions, and expressions involving several independent variables.
What changes is the jet space and the rank of the prolonged group. The method remains the same:
prolong the infinitesimal transformations, solve the invariant equations, and construct invariant
differentiations.


A useful normal form is obtained by choosing a cross-section to the group orbits in a sufficiently
high jet space. The coordinates of the transformed jet at the cross-section are differential
invariants. Once the cross-section is fixed, every other invariant is a function of these normalized
quantities and their invariant derivatives. This normal form is algebraic when the group action is
algebraic and the cross-section is chosen by algebraic equations.


The conclusion is that the invariant theory of differential expressions is structurally parallel to
the invariant theory of algebraic forms. Infinitesimal transformations replace symbolic letters;
prolongation replaces polarization; invariant differentiation replaces transvection. The principal
difference is that the generating process is infinite in order, though finite at each stage and
controlled by finitely many fundamental operations.


\subsection{Additional reductions and formulation}

The prolongation formula is obtained by differentiating the transformed dependent variables and
applying the chain rule. For a first derivative one obtains terms involving derivatives of the
infinitesimal coefficients \(\xi\) and \(\eta\); for higher derivatives the formula is iterated.
Although the expressions become long, they are determined recursively. This recursion is what makes
the invariant equations finite at every prescribed order.

The rank of the prolonged group action may increase with the order of derivatives. At low order the
group may have stabilizers; at sufficiently high order the action may become free on a dense open
set. The number of fundamental differential invariants at each order is then the increase in the
dimension of jet space minus the increase in orbit dimension. This count is the differential
analogue of the dimension count by Hilbert series in algebraic invariant theory.

Invariant differentiation is constructed by finding differential operators that transform
covariantly. If \(\omega_1,\ldots,\omega_s\) is an invariant coframe, the dual total derivative
operators preserve invariants. In coordinates this means solving a linear system whose coefficients
are the transformation coefficients of the independent variables. The determinant of this system
marks the exceptional jets where the construction fails or changes rank.

The method clarifies the role of moving frames even when that terminology is not used. A
cross-section to the group orbits fixes the group parameters as functions of the jet variables.
Substituting these parameters into the transformed derivatives gives normalized differential
invariants. All other invariants are functions of these normalized quantities, provided one stays in
the open set where the cross-section is valid.

Differential syzygies arise because invariant differentiations need not commute. Their commutators
are linear combinations of the invariant differentiations with invariant coefficients. These
coefficients themselves are differential invariants. Consequently, differentiating fundamental
invariants produces relations as well as new invariants. A complete theory must include both the
generating invariants and the commutation relations among invariant differentiations.


\subsection{Further proof details}

At each finite order the differential invariant problem is an ordinary invariant problem on a
finite-dimensional jet space. What is new is the compatibility as the order increases. Prolongation
maps from order \(r+1\) to order \(r\) must carry invariants to invariants, and invariant
differentiations move in the opposite direction. The structure is therefore a graded tower rather
than a single finite algebra.

The arbitrary nature of the differential expression means that the results are not tied to equations
of motion or to a particular geometry. Curvature expressions, differential equations, and
variational integrands are all treated by the same method. One specifies the transformation group,
prolongs it, and solves for invariant functions on jets. The concrete interpretation of the
resulting invariants depends on the problem, not on the general theory.

The finiteness statement should be read locally on the regular part of jet space. Where the rank of
the prolonged action drops, additional invariants may appear or the chosen invariant
differentiations may become singular. This is analogous to algebraic invariant theory, where orbit
dimension may drop on special strata. The general theorem describes the open stratum; special strata
require separate analysis.


\section{13. Invariant variational problems}

\textit{Original publication: Nachrichten der Gesellschaft der Wissenschaften zu Goettingen, 1918, pp. 235--257.}

The problem is to determine the consequences for the Euler equations when a variational integral is
invariant under a continuous group of transformations. Let \[   I=\int f(x,u,\partial u,\ldots)\,dx
\] be an integral whose integrand depends on finitely many derivatives. If the integral is invariant
under a group, then the variational derivatives of \(f\) satisfy identities. Conversely, under
suitable hypotheses such identities imply invariance up to a divergence. This is the content of the
two theorems.


Write the Euler expressions as \[   \psi_i=\frac{\partial f}{\partial u_i}   -\sum_\lambda
\frac{d}{dx_\lambda}\frac{\partial f}{\partial u_{i,\lambda}}
+\sum_{\lambda,\mu}\frac{d^2}{dx_\lambda dx_\mu}        \frac{\partial f}{\partial
u_{i,\lambda\mu}}-\cdots . \] The first variation has the form \[   \delta f=\sum_i \psi_i\delta
u_i+\operatorname{Div} A. \] This formula is the basis of everything that follows. Invariance means
that for every infinitesimal transformation the variation of \(f\) is itself a divergence.


For a finite continuous group depending on \(\rho\) essential parameters, the infinitesimal
transformations have characteristic variations \[   \bar\delta u_i=\sum_{\alpha=1}^\rho a_\alpha
Q_{i\alpha}. \] Substituting in the first variation and using invariance gives \[   \sum_i \psi_i
Q_{i\alpha}=\operatorname{Div} B_\alpha \] for each \(\alpha\). Thus there are \(\rho\) divergence
relations among the Euler expressions. These are conservation laws. They are proper when the
parameters are constant and independent.


This is the first theorem: to every finite continuous symmetry group of the variational integral
there correspond as many independent conservation laws as there are independent group parameters.
The conservation laws are explicitly obtained by inserting the infinitesimal transformations into
the boundary term of the first variation. No integration of the Euler equations is required. The law
is an identity in the functions and their derivatives.


The second theorem concerns infinite groups depending on arbitrary functions and their derivatives.
Suppose the infinitesimal transformations contain arbitrary functions \(p_1(x), \ldots,p_\sigma(x)\)
and their derivatives up to some finite order. After integration by parts, invariance gives
identities of the form \[   \sum_i \psi_i Q_{ij}-\operatorname{Div}\sum_i \psi_i
Q_{ij}^{(1)}+\\cdots=0 \] for \(j=1, \ldots, \sigma\). These are not conservation laws but
dependencies among the Euler expressions themselves and their derivatives. They hold identically,
before any field equations are imposed.


This distinction explains the difference between proper conservation laws and improper ones. For
finite groups, the divergence vanishes on solutions of the Euler equations and gives a conserved
current. For infinite groups, the corresponding divergence is a consequence of identities among the
Euler expressions; it is therefore improper in the sense that it contains no independent
conservation content beyond the identities. This resolves an ambiguity that appears in mechanics and
relativity.


The converse statements are equally important. If \(\rho\) independent divergence relations of the
indicated form are given, then one can reconstruct a finite continuous group under which the
integral is invariant, up to a divergence. If \(\sigma\) independent identities among the Euler
expressions and their derivatives are given, then one obtains an infinite group depending on
\(\sigma\) arbitrary functions. Thus symmetries and identities determine each other.


The proof is algebraic integration by parts. Treat the arbitrary variations and their derivatives as
independent, collect coefficients, and transfer derivatives from the arbitrary functions to the
Euler expressions. Boundary terms are absorbed into divergences. The independence of the arbitrary
functions forces the identities. Conversely, reversing the calculation reconstructs the variation of
the integrand as a divergence.


The theorem applies to classical mechanics, where time translations yield energy conservation and
spatial translations yield momentum conservation. It also applies to the general theory of
relativity. The invariance under arbitrary coordinate transformations is an infinite group.
Therefore the energy relations associated with it are improper conservation laws; they arise from
differential identities among the gravitational field equations. This is the mathematical
explanation of the special status of gravitational energy expressions.


In the language of the calculus of variations, the first theorem associates finite-dimensional
symmetry with conserved divergences; the second associates gauge-type symmetry with differential
identities among Euler equations. The two cases are formally similar but conceptually distinct. The
first yields constants of motion on solutions. The second yields relations that make some field
equations dependent on others.


The final formulation is independent of coordinates. A variational problem is invariant if the
change of the integrand under the prolonged infinitesimal transformation is a total divergence. The
Euler operator annihilates divergences. Therefore applying the Euler operator to the invariance
condition gives the identities; applying the first variation formula gives the conservation laws.
This is the complete mechanism behind the correspondence.


\subsection{Additional reductions and formulation}

The finite-parameter theorem includes the familiar conservation laws as special cases. Time
translation gives the energy integral; translations in space give momentum; rotations give angular
momentum. The theorem does not depend on the mechanical interpretation. It says that every
infinitesimal transformation for which the change of the Lagrangian is a divergence produces a
divergence expression that vanishes on all extremals. The conserved quantity is read directly from
the boundary term of the first variation.

The arbitrary-function theorem is more subtle. If the transformation contains an arbitrary function
and its derivatives, the coefficients of that arbitrary function after integration by parts must
vanish identically. These vanishing expressions are the Noether identities. They imply that not all
Euler equations are independent. In gauge theories this is exactly what one expects: gauge freedom
means that the field variables contain redundant descriptions of the same physical state.

The converse direction is often the most powerful part of the result. Given identities among the
Euler expressions, one can reconstruct transformations with arbitrary functions. This shows that
dependencies among field equations are not accidental; they are the algebraic shadow of an infinite
symmetry group. In particular, if a variational problem has identities of order \(s\), the
corresponding symmetry transformations involve arbitrary functions and derivatives up to a related
order.

Divergence terms must be retained throughout. Two Lagrangians that differ by a divergence have the
same Euler expressions, but their Noether currents differ by an identically conserved term.
Therefore the theorem is naturally formulated modulo divergences. This is not a defect: the
variational problem itself is unchanged by a divergence, and conservation laws are likewise
determined only up to trivial divergences.

The application to general relativity was one of the motivations for the paper. Coordinate
invariance is an infinite group depending on arbitrary functions. The associated energy expressions
cannot be ordinary tensorial conservation laws of the same kind as those coming from finite
symmetries. They are tied to identities among the gravitational equations. The theorem gives the
precise mathematical reason for this phenomenon without relying on physical terminology.

In modern terms the first theorem concerns global symmetries and the second concerns gauge
symmetries. But the paper's formulation is more general than this terminology. It treats finite and
infinite continuous groups uniformly by counting essential constants or arbitrary functions and by
applying integration by parts to the first variation. The result is a structural theorem about
variational complexes, expressed in the calculus notation of the time.


\subsection{Further proof details}

The Euler operator is the reason divergences disappear from the field equations. Applying it to a
divergence gives zero identically. Therefore any invariance statement modulo a divergence becomes,
after applying the Euler operator, an identity among Euler expressions. This simple observation is
the algebraic core of the second theorem and explains why integration by parts is unavoidable in the
proof.

Independence of the conservation laws or identities is measured by the independence of the
infinitesimal transformations. If a linear combination of generators is trivial modulo
transformations that vanish on all fields, the corresponding conservation law is trivial. Thus the
counting of laws requires the group parameters to be essential. Noether states the hypotheses so
that the number of laws is the number of essential parameters or arbitrary functions.

The theorem is local in the independent variables. Boundary conditions are not part of the algebraic
statement. They become relevant only when one integrates the divergence over a domain and interprets
the boundary term as a conserved quantity. The identities themselves hold pointwise as formal
differential identities. This is why the theorem applies equally to mechanics, field theory, and
purely geometric variational problems.


\section{14. Arithmetic theory of algebraic functions of one variable}

\textit{Original publication: Jahresbericht der Deutschen Mathematiker-Vereinigung 28 (1919), pp. 182--203.}

The arithmetic theory of algebraic functions of one variable stands between two older theories: the
theory of algebraic number fields and the geometric theory of algebraic curves. The aim is to
exhibit the precise analogy and also the differences. A function field of one variable over a
constant field behaves like a number field in its ideals and divisors, while its places and
differentials have a geometric interpretation on a curve.


Let \(K/k\) be an algebraic function field of one variable. A prime divisor, or place, is the
valuation corresponding to an irreducible branch of the curve. Divisors are formal sums \[
\mathfrak a=\sum_P n_P P \] with almost all \(n_P=0\). The principal divisors are those of the form
\((f)\), with \(f\in K^\times\). The divisor class group is the quotient of all divisors of degree
zero by principal divisors. It is the function-field analogue of the ideal class group.


The analogy with number fields becomes sharp when one studies integral elements. Choosing a place at
infinity, the functions regular away from infinity form a Dedekind domain. Its non-zero ideals
factor uniquely into prime ideals. The divisor theory of the complete function field is obtained by
adjoining the infinite primes. Thus the arithmetic of ideals and the geometry of points are two
languages for the same structure.


Ramification is measured by the different. If \(L/K\) is a finite separable extension of function
fields, a prime \(P\) of \(K\) decomposes into primes \(Q\) of \(L\), with ramification indices
\(e(Q/P)\) and residue degrees \(f(Q/P)\). The different exponent measures the failure of the trace
pairing to be integral. The sum of these exponents enters the Riemann-Hurwitz formula, just as the
discriminant enters the arithmetic of number fields.


The Riemann-Roch theorem is the central theorem connecting arithmetic and geometry. For a divisor
\(D\), let \(L(D)\) be the vector space of functions \(f\) such that \((f)+D\geq 0\). Then \[
\ell(D)-\ell(K_0-D)=\deg D+1-g, \] where \(K_0\) is a canonical divisor and \(g\) is the genus. In
number-field language this resembles the balance between ideals and differentials; in geometric
language it is the dimension theorem for linear systems.


The theory of differentials gives the most transparent bridge. A differential has orders at all
places, and its divisor is canonical. Residues satisfy the global relation that the sum of all
residues is zero. This is the function-field analogue of the product formula and of trace duality.
Through differentials, the canonical class and the different of an extension acquire arithmetic
meaning.


The constant field must be treated carefully. Extending constants changes the set of rational places
and may split prime divisors, but it does not change the geometric function field after algebraic
closure. Many statements are most natural over an algebraically closed constant field; arithmetic
applications require descent to the original constants. This mirrors the distinction between
absolute and relative ideal theory in number fields.


The paper compares three formulations of the same theory: ideals in an integral domain of functions
regular outside a prescribed divisor, divisors on the complete curve, and valuations of the function
field. Each has advantages. Ideals make unique factorization of ideals immediate. Divisors make
linear equivalence and Riemann-Roch natural. Valuations make extension and ramification transparent.
A complete arithmetic theory must pass freely among the three.


The conclusion is that the function-field case is not a mere analogy to number fields. It is an
exact arithmetic theory in which the geometric model supplies additional tools. Conversely, the
number-field theory can be illuminated by translating its notions into the language of valuations,
divisors, differents, and class groups. This perspective prepares the later abstract construction of
ideal theory in both number and function fields.


\subsection{Additional reductions and formulation}

The divisor attached to a function records zeros and poles with multiplicity. The product formula
says that the total degree of a principal divisor is zero. This is the simplest sign that the theory
is global. Local valuations determine the divisor, but the degree condition ties all places
together. In the number-field analogy, this is the counterpart of the product formula for absolute
values.

The ideal-theoretic and divisor-theoretic languages differ mainly by the treatment of infinity. If
one removes a finite set of places \(S\), the functions regular outside \(S\) form a Dedekind
domain. Ideals in this domain correspond to effective divisors away from \(S\). Adding the places in
\(S\) recovers the complete divisor group. Thus changing the set of infinite places changes the
affine arithmetic but not the complete function field.

The Riemann-Roch theorem also explains why genus is an arithmetic invariant. It measures the defect
in the naive expectation that a divisor of degree \(d\) should impose \(d\) independent conditions.
In genus zero the function field is rational. In higher genus the failure of rational
parametrization is reflected in the correction term involving the canonical divisor. This connects
rationality of one-variable fields with the arithmetic of divisors.

In extensions of function fields the different has two equivalent descriptions. Locally it is
defined by the trace dual of the integral closure of a valuation ring. Globally it is a divisor
whose degree enters Riemann-Hurwitz. The local exponents measure ramification, while the global
degree measures the change of genus. This exact correspondence is one of the strongest reasons for
treating function fields and number fields by the same ideal-theoretic methods.

The constant-field issue appears in the definition of degree of a divisor. If the constant field is
not algebraically closed, a place may have residue degree greater than one. Its contribution to the
divisor degree is weighted by that residue degree. After extending constants, the place may split
into several geometric points. The arithmetic theory over the original field must keep the residue
degrees; the geometric theory over an algebraic closure sees the split points.


\subsection{Further proof details}

The comparison with number fields is especially close for the trace pairing. If \(A\) is the
integral closure of a Dedekind domain in a finite extension, the inverse different consists of
elements whose trace against \(A\) is integral. This definition works in number fields and function
fields alike. In the function-field case, the divisor of the different can also be read
geometrically from ramification of the corresponding curve map.

Linear equivalence of divisors is the function-field form of principal ideal equivalence. Two
divisors differ by a principal divisor exactly when their associated fractional ideals differ by
multiplication by a field element, after choosing the affine ring of functions regular away from
infinity. The class group of divisors is therefore the natural completion of the ideal class group.
Riemann-Roch then gives information about representatives of these classes.

This perspective makes clear why one-variable function fields are the correct analogue of number
fields rather than of arbitrary multivariable function fields. Their local rings at non-singular
points are discrete valuation rings, and their non-zero ideals factor into primes. In higher
dimension prime divisors still exist, but ideals have more complicated primary decompositions. The
one-dimensional case is where the arithmetic analogy is exact.


\section{15. Finiteness of the system of integral invariants of binary forms}

\textit{Original publication: 1919.}

The problem is to prove the finiteness of the system of integral invariants for binary forms. Here
integral means that the invariants are polynomials with integer coefficients in the coefficients of
the ground forms. The usual proofs over fields do not immediately give integral generators, since
division by integers may occur. The task is to avoid such division or to control it so that the
final generators remain integral.


Let a binary form of degree \(n\) be written \[   f=a_0x^n+a_1x^{n-1}y+\cdots+a_ny^n. \] The group
\(SL_2\) acts on the variables and hence on the coefficients. An integral invariant is a polynomial
\(I(a_0, \ldots,a_n)\in\mathbb Z[a_0, \ldots,a_n]\) fixed by this action. For several binary forms
one takes the diagonal action on all coefficient sets. The invariant ring over \(\mathbb Z\) is the
object to be generated.


Over a field of characteristic zero, Hilbert's theorem gives finite generation. But the proof uses
averaging or differential operators that may introduce rational coefficients. One may clear
denominators for any one invariant, but this does not by itself produce a finite integral generating
system. The integral theorem requires showing that all denominators needed in the field proof can be
absorbed into finitely many integral forms.


The method is to compare the invariant ring over \(\mathbb Z\) with its extensions to \(\mathbb Q\)
and to the residue fields modulo primes. Let \(R_\mathbb Z\) be the integral invariant ring and
\(R_\mathbb Q\) its rational extension. Choose rational generators of \(R_\mathbb Q\), clear
denominators, and obtain finitely many integral invariants. They generate a subring \(A\subset
R_\mathbb Z\) such that \(R_\mathbb Z/A\) is torsion. The remaining question is to control this
torsion uniformly.


Reduction modulo a prime \(p\) detects the possible missing integral invariants. If an invariant is
not generated by \(A\), its residue modulo some prime gives a new invariant in characteristic \(p\).
Therefore it is enough to prove finite generation in every characteristic in a manner compatible
with lifting. For binary forms the representation theory supplies enough operators to reduce
high-degree invariants to lower-degree ones, even modulo \(p\), after separating the exceptional
primes.


The symbolic method again plays a role. Transvectants of binary forms can be written with integral
coefficients if the binomial factors are normalized correctly. The apparent divisions in classical
transvectants are artifacts of convention. By choosing integral symbolic operators, one obtains
covariants and invariants whose coefficients lie in \(\mathbb Z\). Products and contractions of
these integral covariants remain integral.


A key lemma states that, beyond a computable degree, every invariant is congruent modulo the ideal
generated by positive-degree invariants of lower degree to an invariant of a restricted normal form.
The number of such normal forms is finite in each degree range. Combined with the rational
finite-generation theorem, this gives a bound beyond which no new integral generator is required.
The argument is a graded version of the usual ascent from quotient modules.


The theorem then follows from the noetherian property of the polynomial ring over \(\mathbb Z\). The
invariant ring is a graded subring, and the reductions show that it is integral over the finitely
generated subring \(A\) up to bounded torsion. Since only finitely many primes and finitely many
degrees remain exceptional, adjoining finitely many further integral invariants completes the
generating system.


The conclusion is that the full ring of integral invariants of any finite system of binary forms is
finitely generated over \(\mathbb Z\). This is stronger than finite generation over \(\mathbb Q\):
it asserts that a single finite list of invariant polynomials with integer coefficients generates
all integral invariants without introducing denominators. The result is an integral refinement of
Hilbert's finiteness theorem.


\subsection{Additional reductions and formulation}

One begins with a graded invariant ring. The degree-zero part is \(\mathbb Z\), and the
positive-degree part is an ideal. If a finite set of homogeneous invariants generates this ideal up
to a given degree and if every higher-degree invariant is reducible modulo lower degrees, the whole
ring is generated by the finite set. The proof is the standard induction on degree, but the
difficulty is to make the reduction integral.

The passage from rational to integral invariants can be organized prime by prime. For a prime \(p\),
reduce the coefficient ring modulo \(p\) and consider invariants over \(\mathbb F_p\). If the
reductions of the chosen integral generators generate the invariant ring in characteristic \(p\) up
to the required bound, then no new generator with denominator divisible by \(p\) is needed. Only
finitely many primes can fail the chosen finite checks.

Binary forms are special because the \(SL_2\)-operation is controlled by two fundamental
differential operators, raising and lowering the weight. In characteristic zero this gives complete
reducibility. Integrally, the operators must be normalized to avoid denominators. The resulting
congruence arguments show that high-weight terms can be reduced to transvectants and products
already in the generated ring.

The theorem does not assert that the integral generating system is minimal. Minimality over
\(\mathbb Z\) is not stable under reduction modulo primes, and generators that are redundant over
\(\mathbb Q\) may be needed integrally to account for congruence phenomena. The statement is finite
generation: there exists one finite list that suffices for all integral invariants. This is the
right analogue of Hilbert's theorem over the integers.

The result is important because it shows that invariant theory over the integers is not merely
rational invariant theory with denominators cleared case by case. It is a noetherian theory in its
own right. Integral invariants form a finitely generated algebra, and arithmetic questions about
congruences of invariants can be studied inside that algebra without returning to infinitely many
separate denominators.


\subsection{Further proof details}

The proof also requires that products of invariants behave well under reduction. If an invariant
reduces modulo \(p\) to a product of lower-degree invariants, then after lifting the factors one
obtains the original invariant modulo \(p\) and lower-degree terms. Repeating this process and using
induction on degree reduces the integral problem to finitely many congruence classes. This is the
arithmetic form of the reduction argument.

The exceptional primes are unavoidable because representation theory of \(SL_2\) changes in
characteristic \(p\). Operators that are invertible over \(\mathbb Q\) may acquire kernels modulo
\(p\), and new invariants may appear in small characteristic. The theorem allows this: one adds
finitely many integral lifts to handle the exceptional primes. What it excludes is the need for
infinitely many new generators as the degree grows.

The final integral generating system can be chosen homogeneous. This follows from the grading by
degree in the coefficients. If a non-homogeneous invariant is present, each homogeneous component is
invariant as well, because the group action preserves degree. Homogeneous generators are therefore
sufficient and make the induction on degree possible.


\section{17. Modules in non-commutative domains, especially from differential and difference expressions}

\textit{Original publication: Mathematische Zeitschrift 8 (1920), pp. 1--35.}

The theory of modules over commutative polynomial rings had already shown the power of
ideal-theoretic methods. The purpose here is to extend part of that theory to non-commutative
domains that arise naturally from differential and difference operators. In such domains
multiplication depends on commutation rules, and left ideals need not be right ideals. One must
therefore formulate every statement with attention to side.


The basic examples are operator rings. In a differential operator ring one has a coefficient field
\(K\) and an operator \(D\) satisfying \[   D a=aD+D(a) \] for \(a\in K\). In a difference operator
ring one has an automorphism \(S\) and a rule \[   S a=a^S S. \] Both rings are non-commutative
unless the derivation or automorphism is trivial. Nevertheless they admit division algorithms in
suitable cases and therefore possess ideal-theoretic structure.


A left module over such a ring represents a system of linear differential or difference equations.
Generators correspond to unknown functions; relations correspond to equations. Transforming the
system by elementary operations is the same as replacing the module presentation by an equivalent
presentation. Thus module theory gives an invariant language for systems of operators.


The first question is the chain condition. A domain is left-Noetherian if every ascending chain of
left ideals stabilizes. For operator rings of the above type, the proof uses an order filtration.
The leading term of an operator is commutative enough to be compared by degree. An ascending chain
of left ideals gives an ascending chain of leading coefficient ideals or leading monomial sets,
which stabilizes. One then descends on order to show that the original chain stabilizes.


Because left and right multiplication differ, one must also study right ideals separately. In many
operator rings the same filtration proves both left and right chain conditions, but the generators
and reductions are not identical. The distinction becomes important when forming quotient modules
and when asking for common right or left multiples of operators. The classical least common multiple
of differential operators is a one-sided concept.


The elementary divisor theorem does not carry over unchanged. Over a principal ideal domain every
finitely generated module decomposes into cyclic modules. In non-commutative domains, cyclic
decompositions may fail or may depend on side. What remains is a theory of presentations, syzygies,
and reductions to normal form under elementary row and column operations that respect the
non-commutative multiplication rules.


For differential equations, this means that a system can be simplified by eliminating unknowns and
equations, but the resulting single equation, when it exists, is not canonical in the same way as in
the commutative case. The annihilator of a module and the common multiples of the operators encode
the invariant content. The module itself is the correct object; a chosen scalar equation is only one
representation of it.


The paper develops divisibility notions adapted to non-commutative rings. If \(a=bc\), then \(b\) is
a left divisor and \(c\) a right divisor. Greatest common divisors and least common multiples must
be qualified accordingly. Euclidean arguments, when valid, produce one-sided greatest common
divisors. These are sufficient for many transformations of differential and difference systems.


The conclusion is that a substantial part of ideal and module theory survives in non-commutative
operator domains once one respects the side conditions and uses filtrations. This opens the way to
treating differential and difference equations by algebraic methods analogous to the ideal theory of
polynomial equations. Later non-commutative algebra will make these ideas systematic, but the
essential module viewpoint is already present here.


\subsection{Additional reductions and formulation}

The order filtration is the central technical tool. Assign to an operator its order in \(D\) or
\(S\). The leading coefficient of a product is the product of the leading coefficients, with the
appropriate automorphism in the difference case. Therefore leading terms behave sufficiently like
commutative monomials to permit induction. Lower-order commutator terms do not destroy the argument
because they are handled after the leading term has been reduced.

Syzygies must also be one-sided. A relation among left module generators has coefficients acting
from the left. If the same symbols are moved to the right, the relation changes. This is not a
notational inconvenience but an algebraic fact. For systems of differential equations it corresponds
to the distinction between combining equations by applying operators to them and combining unknown
functions by operator substitutions.

A finitely generated left module over a left-Noetherian operator ring has a finite presentation.
This means that the relations among a finite set of equations are themselves finitely generated. In
analytic language, compatibility conditions of a linear differential system can be generated from
finitely many identities. The algebraic theorem does not solve the equations, but it gives a finite
algebraic description of their formal consequences.

The comparison with commutative elementary divisors shows both what survives and what fails. Torsion
modules still make sense: an element is torsion if some non-zero operator annihilates it. But
annihilators are left ideals, and their factorization is not controlled by commutative prime ideals.
Decomposition theory must therefore be replaced by ideal chains, filtrations, and one-sided
divisibility.

The examples from difference equations are parallel to the differential ones but reveal the role of
automorphisms more clearly. The rule \(Sa=a^S S\) means that coefficients are shifted when an
operator passes them. Common multiples correspond to bringing equations to a common shift range. The
algebraic operations are the same as for differential operators after replacing derivation by
automorphism in the commutation rule.

The paper's broader message is methodological. Non-commutative algebra should not be treated as a
collection of exceptional calculi. Operator domains have ideals, modules, chain conditions, quotient
modules, and presentations, provided the definitions are made on the correct side. This is the point
of departure for a structural theory of non-commutative rings analogous to the ideal theory of
commutative domains.


\subsection{Further proof details}

The filtration also gives an associated graded ring. In the differential case the associated graded
ring is often commutative or at least simpler than the original operator ring. Ideals and modules
can then be studied by passing to leading terms in the associated graded ring and lifting reductions
back to the filtered ring. This is the method behind the chain-condition proofs.

For systems of equations, finite presentation has a concrete computational meaning. Starting from
finitely many equations, there are finitely many generating compatibility conditions; compatibility
conditions among those are again finitely generated, and so on in each finite stage allowed by the
ring's homological behavior. The paper does not use modern homological language, but it already
treats equations through their modules of relations.

Non-commutativity forces care with localization. In a commutative domain one may invert a
multiplicative set without choosing sides. In an operator domain, denominators must satisfy
compatibility conditions if a ring of fractions is to exist. The differential and difference
examples often have enough Ore-type behavior for the necessary fractions, but this is a theorem, not
a convention. The module theory must respect this limitation.


\end{document}
